\section{Анализ предметной области}

\subsection{Характеристика предметной области}

Написание книги -- это трудоемкий процесс. Помимо работы над основным текстом, автору приходится работать с множеством вспомогательного материала, такого как заметки по сюжету, описание персонажей, проработка нескольких сюжетных линий одновременно. Процесс написания редко идёт строго по порядку. Автор может переключаться в процессе написания между главами, менять фрагменты местами, объединять или разделять их. Придерживаться четкой структуры написания текста без вспомогательных средств становится всё труднее.

На практике, чтобы не запутаться в материалах, авторы, как правило, приходят к одному и тому же решению -- делят книгу на части. Сначала идут сюжетные арки -- крупные блоки, объединённые общей линией, -- а внутри них уже главы. Такая иерархия позволяет видеть произведение целиком и не терять связь между разными его частями.

Цифровая среда дала авторам новые инструменты для работы, но саму задачу организации текста никто не отменял. Когда под рукой десятки файлов, заметок и черновиков, потребность в чёткой структуре только растёт. Помимо творческой работы, автору приходится заниматься и учётом результатов: следить за количеством слов и символов по главам и книге в целом, оценивать темп работы и планировать сроки завершения.

\subsection{Проблематика традиционного подхода к организации писательской работы}

В настоящее время большинство авторов применяют для работы текстовые редакторы общего назначения в сочетании с файловой системой. Главы сохраняются как отдельные документы, которые группируются по папкам, соответствующим аркам книги. Такой подход прост, но имеет несколько ограничений.

Хранение глав в виде набора файлов не даёт единого представления о структуре книги. Навигация между главами ведётся через файловый менеджер. При большом количестве глав это превращается в рутину.

Подсчёт статистики по книге тоже выполняется вручную. Автор вынужден суммировать показатели из отдельных файлов либо использовать внешние сервисы. Процедура несложная, но её регулярное повторение отнимает время.

Импорт материалов из внешних источников -- черновиков в DOCX или PDF -- делается через буфер обмена и требует ручного форматирования. Экспорт готового текста для редактора также предполагает сборку глав в один документ.

В традиционных редакторах сохранение текста выполняется вручную. Если автор забыл сохранить документ и произошёл сбой программы или отключение питания, часть работы может быть потеряна. Автосохранение, реализованное на стороне сервера, позволяет избежать таких ситуаций.

Наконец, доступ к языковым моделям реализован вне рабочей среды автора. Чтобы обратиться к ИИ, нужно открыть сторонний сервис, скопировать туда текст, описать задачу. Это отвлекает и прерывает творческий процесс.

\subsection{Обоснование необходимости автоматизации}

Эти проблемы показывают, что для крупных проектов нужен другой инструмент, а именно веб-приложение, в котором автор работает с книгой в одном месте, не переключаясь между программами.

Заменой файловой системе может служить древовидная структура: арки раскрываются, внутри -- список глав. Автор видит книгу целиком и быстро переходит к нужному разделу. Все данные лежат на сервере, а значит, доступны с любого устройства.

Встроенный редактор позволяет работать с форматированием без обращения к внешним программам. Автосохранение на сервере защищает от потери текста при сбоях. Приложение само считает слова и символы по главам -- не надо делать это вручную.

Отдельного внимания заслуживает работа с внешними форматами. Возможность загрузки черновиков, подготовленных в DOCX или сохранённых в PDF, без потери содержимого позволяет автору переносить наработки из других инструментов. Готовые главы, напротив, необходимо экспортировать в DOCX с сохранением форматирования -- это стандартный формат для передачи текста редакторам и издательствам.

В числе прочего в приложении целесообразно предусмотреть встроенного ИИ-ассистента, работающего через API языковых моделей. Его вызов должен осуществляться непосредственно из интерфейса системы, без необходимости перехода на сторонние ресурсы. Предполагается два режима работы. В контекстном режиме ассистент опирается на текст текущей главы и фрагменты других глав, что позволяет давать ответы, учитывающие сюжет и персонажей. В общем режиме он обрабатывает запросы, не связанные с содержанием книги: помогает подобрать формулировку, предлагает литературный приём, отвечает на вопросы общего характера. Переключение между режимами выполняется в самом приложении,благодаря чему автору не требуется копировать фрагменты текста в сторонние сервисы.

Когда все эти функции собраны в одном приложении, автор тратит меньше времени на технические операции и может сосредоточиться на работе над текстом.

\subsection{Анализ существующих систем управления авторским контентом}

На рынке представлен ряд программных продуктов, которые частично решают задачи организации писательской работы. Ниже рассмотрены наиболее распространённые решения, их функциональные возможности и ограничения применительно к задаче управления авторским контентом с иерархической структурой.

\subsubsection{Microsoft Word}

Microsoft Word является наиболее распространённым текстовым редактором, который используют в том числе и авторы книг. Его возможности включают развитые средства форматирования текста, проверку орфографии и грамматики, режим рецензирования, удобный для работы с редакторами, облачное автосохранение при использовании OneDrive, а также кроссплатформенность -- приложение доступно на Windows, macOS и в веб-версии.

Однако Word не рассчитан на ведение крупных проектов. В нём нет иерархии «книга--арка--глава»: текст либо хранится в одном файле, где трудно ориентироваться, либо разбивается на отдельные файлы, из-за чего теряется целостное представление о книге. Нет средств для отслеживания прогресса по главам, а ИИ-ассистент Copilot работает с общими задачами оформления и не учитывает сюжетный контекст произведения.

\subsubsection{Scrivener}

Scrivener является одним из наиболее функциональных десктопных инструментов для писателей. Программа включает разбиение проекта на части, главы и сцены с возможностью изменения порядка перетаскиванием, хранение исследовательских заметок, карточек персонажей и локаций, визуализацию структуры произведения через режим «пробковой доски», отслеживание прогресса с установкой целевого количества слов, а также встроенный компилятор для сборки рукописи и экспорта в DOCX, ePub, PDF.

Однако Scrivener имеет ряд ограничений, которые делают его не вполне подходящим для решаемой задачи. Это исключительно десктопное приложение, что не позволяет автору работать из любого места и с любого устройства. Синхронизация между устройствами требует ручного переноса файлов или настройки облачных сервисов. Для российских пользователей приобретение лицензии затруднено из-за ограничений на оплату зарубежного программного обеспечения. Кроме того, интеграция с ИИ-ассистентами в программе не предусмотрена, что не позволяет автору получать интеллектуальную помощь в процессе работы. Эти ограничения делают Scrivener непригодным для создания веб-ориентированного решения с интеграцией ИИ.

\subsubsection{yWriter}

yWriter представляет собой бесплатную программу для структурирования книг, разработанную специально для писателей, которым необходимо организовать работу над крупными проектами. Приложение позволяет разбивать книгу на главы и сцены с автоматическим подсчётом статистики по каждой из них, что избавляет автора от необходимости вручную отслеживать объём написанного текста. Также в yWriter реализовано ведение учёта персонажей, локаций и заметок -- вся вспомогательная информация хранится в одном месте и привязана к конкретному проекту. Пользователь может устанавливать целевые показатели по количеству слов и отслеживать прогресс выполнения, а также экспортировать готовый текст в формат DOCX для передачи редакторам или издательствам.

Несмотря на свою функциональность, yWriter имеет существенные ограничения для решаемой задачи. Программа работает только на операционной системе Windows -- веб-версия и мобильные клиенты отсутствуют, что делает невозможным доступ к проекту с других устройств и исключает возможность совместной работы. Интерфейс приложения заметно устарел и не обновлялся в течение нескольких лет, что негативно сказывается на удобстве использования и восприятии информации. Кроме того, в программе не реализованы функции импорта черновиков, что затрудняет перенос наработок из других инструментов, а также отсутствует интеграция с ИИ-ассистентами, которая могла бы помочь автору в творческом процессе. Эти ограничения делают yWriter не вполне подходящим выбором для современного писателя, работающего с разных устройств и использующего возможности искусственного интеллекта.

\subsubsection{Reedsy Studio}

Reedsy Studio представляет собой веб-сервис, предоставляющий авторам инструменты для полного цикла работы над книгой -- от черновика до публикации. Сервис работает по условно-бесплатной модели, разделяя функции на бесплатные и платные.

Бесплатные возможности включают полноценную работу через браузер без установки дополнительного программного обеспечения, иерархическую структуру рукописи с разбивкой на части и главы, базовое отслеживание прогресса с 30-дневной историей версий, совместную работу в реальном времени, а также автоматическую вёрстку и экспорт в ePub и PDF. Платные функции включают расширенную статистику, неограниченную историю версий и инструменты для планирования сюжета и персонажей.

Несмотря на широкий функционал, основным ограничением Reedsy Studio для решаемой задачи является отсутствие встроенного ИИ-ассистента. Кроме того, использование зарубежного облачного сервиса создаёт для российских пользователей риски, связанные с возможными перебоями в доступе и сложностями с оплатой платных тарифов. Эти недостатки делают Reedsy Studio неоптимальным выбором для разрабатываемой системы.

\subsubsection{Writer Journal}

Writer Journal -- это мобильное приложение, предназначенное для написания книг на смартфонах и планшетах под управлением Android. Оно оснащено WYSIWYG-редактором с поддержкой базового форматирования текста, позволяет импортировать документы в формате DOCX, а также экспортировать готовый текст в несколько распространённых форматов.

Однако для решения поставленной задачи Writer Journal подходит плохо. Приложение является сугубо мобильным -- полноценная веб-версия для работы с компьютером отсутствует, что не позволяет автору удобно редактировать большие объёмы текста за рабочим столом. Иерархическая структура книги не поддерживается: весь текст организован в виде набора файлов внутри одной папки, без возможности деления на сюжетные арки и главы. Автор не может видеть структуру произведения целиком и быстро переключаться между его частями. Кроме того, в приложении отсутствует встроенный интеллектуальный ассистент, который мог бы помогать с генерацией идей, редактированием или ответами на вопросы по тексту. Таким образом, Writer Journal не подходит для решения поставленной задачи.

\subsubsection{Вывод по анализу аналогов}

Проведённый анализ показывает, что ни одно из рассмотренных решений не обеспечивает полного набора требуемых функций: веб-доступ, иерархическую структуру «книга--арка--глава» в сочетании со встроенным ИИ-ассистентом. Десктопные приложения проигрывают в кроссплатформенности, веб-сервисы либо не поддерживают нужную структуру, либо не имеют ИИ-интеграции, а мобильные приложения не заменяют полноценную рабочую среду. Это подтверждает целесообразность разработки собственного веб-приложения, ориентированного на поддержку творческого процесса написания.
